% TOP_PROBLEM: {"id": 2306807, "title": "Research Problems in Function Theory — Problem 6.7′", "category_id": 8, "category": {"id": 8, "name": "analysis", "display_name": "Analysis", "description": "Limits, continuity, calculus, and function theory.", "slug": "analysis", "order_index": 8, "created_at": "2026-07-31T15:26:25.671Z"}, "status": "open"}
% FIRST_POSTED: null
% ENTRY_KIND: statement_only
% Imported from: batch04.tex; UnsolvedMath ID 2306807
\documentclass[11pt]{article}
\usepackage[margin=25mm]{geometry}
\usepackage{fontspec}
\setmainfont{Latin Modern Roman}
\usepackage{amsmath,amssymb,mathtools}
\usepackage{unicode-math}
\setmathfont{Latin Modern Math}
\usepackage{xurl,hyperref}
\hypersetup{colorlinks=true,urlcolor=blue}
\setcounter{secnumdepth}{0}
\setlength{\parindent}{0pt}
\setlength{\parskip}{5pt}
\setlength{\emergencystretch}{3em}
\newcommand{\sref}[2]{\hyperlink{source:#1:#2}{[S#2]}}
\newcommand{\cataloguescope}[1]{\paragraph{Recorded target.}#1}
\begin{document}
% UnsolvedMath ID 2306807; source catalogue code AMR-022-6807
\setcounter{section}{2306806}
\section{Research Problems in Function Theory — Problem 6.7$\prime$}\label{2306807}
{\small\sffamily UnsolvedMath ID 2306807 · Analysis}
\subsection{Definitions and mathematical statement}

\paragraph{Shared notation.}$\mathbb N=\{1,2,\ldots\}$, and $\mathbb Z,\mathbb Q,\mathbb R,\mathbb C$ denote the integers, rationals, reals and complex numbers. Any other conventions are specified locally.


Write $\mathbb D=\{z\in\mathbb C:|z|<1\}$ and $\mathbb T=\partial\mathbb D$. For $p>0$ and a holomorphic $f$ on $\mathbb D$, let $n_f(w)$ count the preimages of $w$ with multiplicity. The function is areally mean $p$-valent when
\[\int_0^{2\pi}\int_0^R n_f(\rho e^{i\theta})\rho\,d\rho\,d\theta\leq\pi pR^2\qquad(R>0).\]
Writing $f(z)=\sum_{j=0}^{\infty}a_jz^j$, put $\mu_p(f)=\max\{|a_j|:j\in\mathbb Z,\ 0\leq j\leq p\}$, and let $S(p)$ be the mean $p$-valent functions with $\mu_p(f)=1$.
Here bounded means $\sup_{z\in\mathbb D}|f(z)|<\infty$. The relation $u_n=o(n^{-1/2})$ means $n^{1/2}u_n\to0$, while $u_n=O(n^{-1/2-\varepsilon})$ means $n^{1/2+\varepsilon}|u_n|$ is eventually bounded.
\paragraph{Statement recorded in the supplied source.}
Record the sharp decay assertion stated in the source: for every $p>0$ and bounded $f\in S(p)$,
\[|a_n|=o(n^{-1/2}),\]
and no fixed improvement of the exponent is valid throughout the class: for every $p>0$ and $\varepsilon>0$ there is a bounded $f\in S(p)$ whose coefficients are not $O(n^{-1/2-\varepsilon})$. The source presents this as a consequence of its preceding counterexample, rather than as a remaining conjecture. \sref{2306807}{2}
\subsection{Short English statement}
Bounded mean-valent functions have little-o square-root coefficient decay, but no uniform improvement to a larger power exponent.
\subsection{Sources}
\begin{description}
\item[\hypertarget{source:2306807:1}{S1}] ULAM.AI and the UnsolvedMath Contributors, UnsolvedMath: A Curated Collection of Open Mathematics Problems, record 2306807 (AMR-022-6807). CC BY 4.0. \url{https://huggingface.co/datasets/ulamai/UnsolvedMath}.
\item[\hypertarget{source:2306807:2}{S2}] W. K. Hayman and E. F. Lingham, Research Problems in Function Theory (2018), Problem 6.7 prime, complete problem and update, complete preceding sparse-coefficient counterexample and normalization. \url{https://arxiv.org/abs/1809.07200}.
\end{description}
\label{end:2306807}
\end{document}
